\chapter{Introducción}
\label{cap:introduccion}

El presente proyecto aborda el diseño, implementación y validación de un sistema de monitorización y control de confort térmico para el hogar que opera de forma completamente autónoma en una red local privada, eliminando cualquier dependencia de servicios en la nube para su funcionamiento. El sistema integra un enchufe inteligente, un sensor de temperatura y humedad y un ventilador de techo, todos ellos controlados mediante inteligencia artificial local que interpreta comandos en lenguaje natural y toma decisiones de confort sin enviar datos al exterior.

Este trabajo nace de una observación sencilla: los hogares inteligentes de hoy no son tan inteligentes como parecen. La mayoría de los dispositivos domóticos dependen de servidores externos para funcionar, de modo que cuando falla la conexión a Internet, o cuando el fabricante decide dejar de dar soporte, el usuario se queda sin poder encender una luz o regular el termostato.

El problema no es menor. Cada vez que un sensor de temperatura envía datos a la nube, o que un asistente de voz procesa una orden en servidores remotos, la privacidad del usuario queda expuesta. Los patrones de comportamiento doméstico (a qué hora se levanta, qué temperatura prefiere para dormir, qué días está en casa) viajan a servidores de terceros, a menudo sin que el usuario sea consciente ni haya dado un consentimiento informado.

El proyecto que se presenta en esta memoria, \textbf{Luxe Core AI}, implementa esta idea: en lugar de delegar el procesamiento a la nube, los modelos de lenguaje se ejecutan en un servidor doméstico, y los datos de los sensores nunca abandonan la red local.

Esta decisión de diseño tiene consecuencias prácticas. La latencia de respuesta se reduce de segundos a milisegundos para la mayoría de los comandos. El sistema sigue funcionando aunque el router pierda la conexión con el exterior. Y, sobre todo, los datos del hogar pertenecen exclusivamente a quien vive en él.

Una particularidad del proyecto es que integra dispositivos que no fueron diseñados para interoperar. El proyecto incluye un enchufe inteligente Tuya controlado sin pasar por la nube del fabricante (que a su vez alimenta un humidificador ultrasónico de 24~V, un dispositivo originalmente ``tonto'' que solo necesita corriente y que, al conectarlo al enchufe, adquiere control remoto), un sensor de temperatura Xiaomi leído mediante Bluetooth Low Energy sin necesidad de su aplicación oficial, y un ventilador de techo de 35~\euro{} adquirido en AliExpress cuyo protocolo de comunicación fue descubierto mediante ingeniería inversa. Todos estos dispositivos, que individualmente son ``tontos'', adquieren capacidad de razonamiento al ser orquestados por la capa de inteligencia artificial. El resultado es que la interoperabilidad no depende de estándares cerrados ni de la buena voluntad de los fabricantes, sino de una capa de software que entiende y coordina hardware heterogéneo.

El desarrollo ha seguido dos fases claramente diferenciadas. Durante la primera, se utilizó la infraestructura de la Universidad de Córdoba (servidores con GPU RTX accesibles por VPN) para prototipar la arquitectura de agentes y validar que el enfoque era viable. En la segunda fase, todo el sistema migró a una estación de trabajo doméstica, eliminando cualquier dependencia externa. Esta evolución no fue un retroceso, sino el paso natural hacia el objetivo central del proyecto: comprobar que el control del entorno doméstico mediante IA puede funcionar sin depender de infraestructuras externas.

\section{Motivación Personal}

Más allá de las consideraciones técnicas, este proyecto responde a una motivación personal: aprender a trabajar con tecnologías de \textit{Edge AI} en un entorno real, con hardware real, y comprobar por mí mismo hasta dónde pueden llegar los modelos de lenguaje pequeños ejecutados en local.

Cuando empecé el TFG, los modelos de lenguaje grandes (más de 30~B parámetros) copaban todos los titulares, pero requerían GPUs de varios miles de euros que estaban fuera de mi alcance. En paralelo, modelos mucho más ligeros (como \texttt{qwen2.5-coder:1.5b} o \texttt{bge-m3}) estaban demostrando que se podía hacer mucho con muy poco: clasificación de comandos, análisis semántico, generación de texto coherente, todo en menos de 500~MB de RAM. Quería comprobar si esos modelos eran suficientemente buenos para un caso de uso real, con las limitaciones de un servidor doméstico. El proceso de evaluación, detallado en la Sección~\ref{sec:migracion_qwen}, condujo a una arquitectura mixta: Qwen2.5 1.5B como clasificador determinista y Mistral 7B como razonador conversacional, ambos ejecutándose en CPU con cuantización optimizada.

Los resultados han superado mis expectativas. El clasificador 1.5B responde en 300--500~ms con JSON determinista, el razonador Mistral 7B conversa con naturalidad en 27--60~s, y el sistema completo funciona 24/7 sin intervención. Esta experiencia me ha confirmado que los modelos pequeños y locales son una opción viable para aplicaciones con recursos limitados, y que su capacidad sigue mejorando.

A nivel de documentación, esta memoria se ha redactado con \LaTeX{} y se gestiona con el mismo control de versiones (Git) que el código del proyecto. Esto permite mantener la trazabilidad entre las decisiones de diseño, el código que las implementa y la documentación que las describe.

Luxe Core AI es, en definitiva, un caso práctico de cómo la inteligencia artificial puede aplicarse al entorno doméstico sin renunciar a la privacidad ni depender de infraestructuras externas.

Cabe señalar que la estructura de esta memoria no sigue de forma literal el orden de la guía de desarrollo oficial del GII~\cite{uco_tfg_guide}, pues se ha optado por una organización temática que facilita la lectura del proceso completo. No obstante, toda la información exigida por el reglamento (introducción, objetivos, antecedentes, análisis del problema, diseño, conclusiones, bibliografía y manuales) está presente y es fácilmente localizable a lo largo de los capítulos y anexos.

